This report documents the design, implementation and evaluation of a prototype
client--server application that provides application-layer confidentiality for
messages and bulk files using three ciphers of increasing sophistication: the
Caesar substitution cipher, the Playfair digraph cipher and a byte-wise
Simplified Data Encryption Standard (S-DES). The prototype is written in Python
and consists of four components: a reusable cryptographic module
(\texttt{ciphers.py}), a TCP server (\texttt{server.py}), an interactive TCP
client (\texttt{client.py}) and a sample-data generator
(\texttt{make\_files.py}). The client encrypts user input or a 1~MB binary file,
prefixes it with a four-byte length field and a JSON header describing the mode,
algorithm and key, and transmits the frame over a TCP socket to the server on
port~9000. The server reconstructs the frame, decrypts the payload, displays both
ciphertext and recovered plaintext, and answers with an encrypted reply or an
S-DES-encrypted 10~KB file. Functional testing confirmed loss-free round trips
for all three ciphers: the 1~MB file recovered by the server was byte-identical
to the original, as was the 10~KB file returned to the client. Measured
throughput of the educational S-DES implementation was approximately
0.13~MB/s (7.50~s to encrypt and 7.52~s to decrypt 1~MB), which quantifies the
cost of a pure-Python, block-at-a-time design. The testing chapter also records
the limits observed in practice: Caesar offers 25 effective keys, S-DES offers a
10-bit key space of 1{,}024 keys, encryption is performed in Electronic Codebook
mode so identical plaintext bytes produce identical ciphertext bytes, no
integrity or authentication mechanism is present, and the session key travels in
cleartext inside the JSON header. The prototype therefore achieves obfuscation
suitable for teaching rather than real confidentiality.
